% =============================
% Tractable Circuit Zoo - Transformation Support
% Auto-generated from database.json
% =============================
\documentclass[11pt]{article}
\usepackage[margin=1in]{geometry}
\usepackage{amsmath, amssymb, amsthm}
\usepackage{mathtools}
\usepackage{xparse}
\usepackage{hyperref}
\usepackage{natbib}
\hypersetup{colorlinks=true, linkcolor=blue, citecolor=blue, urlcolor=blue}



\NewDocumentCommand{\langref}{m g}{\textbf{#1\IfNoValueF{#2}{#2}}}
\NewDocumentCommand{\langfam}{m m g}{\textbf{#1$_{#2}$\IfNoValueF{#3}{#3}}}
\newcommand{\thislang}{\ensuremath{\mathcal{L}}}
\newcommand{\poly}{polynomial}
\newcommand{\nopolyunknownquasi}{not polynomial, quasipolynomial unknown}
\newcommand{\nopolyquasi}{not polynomial, quasipolynomial}
\newcommand{\unknownpolyquasi}{polynomial unknown, quasipolynomial}
\newcommand{\unknownboth}{unknown}
\newcommand{\noquasi}{not quasipolynomial}
\newcommand{\CO}{CO}
\newcommand{\VA}{VA}
\newcommand{\CE}{CE}
\newcommand{\IM}{IM}
\newcommand{\EQ}{EQ}
\newcommand{\SE}{SE}
\newcommand{\CT}{CT}
\newcommand{\ME}{ME}
\newcommand{\CD}{CD}
\newcommand{\FO}{FO}
\newcommand{\SFO}{SFO}
\newcommand{\NOTC}{$\neg$C}
\newcommand{\ANDC}{$\wedge$C}
\newcommand{\ANDBC}{$\wedge$BC}
\newcommand{\ORC}{$\vee$C}
\newcommand{\ORBC}{$\vee$BC}
\newcommand{\compilespoly}[2]{#1 compiles to #2 with polynomial blowup}
\newcommand{\compilesquasi}[2]{#1 compiles to #2 with quasipolynomial blowup}
\newcommand{\nocompilespoly}[2]{#1 does not compile to #2 with polynomial blowup}
\newcommand{\nocompilesquasi}[2]{#1 does not compile to #2 with quasipolynomial blowup}
\newcommand{\supportspoly}[2]{#1 supports #2 in polynomial time}
\newcommand{\supportsquasi}[2]{#1 supports #2 in quasipolynomial time}
\newcommand{\nosupportspoly}[2]{#2 is not supported by #1 in polynomial time}
\newcommand{\nosupportsquasi}[2]{#2 is not supported by #1 in quasipolynomial time}
\newcommand{\isin}[1]{#1}
\newcommand{\shortname}[1]{\textbf{Short name:} #1\par}
\newcommand{\fullname}[1]{\textbf{Full name:} #1\par}
\newcommand{\source}[1]{\textbf{Source:} #1\par}
\newcommand{\target}[1]{\textbf{Target:} #1\par}
\newcommand{\claimlanguage}[1]{\textbf{Language:} #1\par}
\newcommand{\operation}[1]{\textbf{Operation:} #1\par}
\newcommand{\status}[1]{\textbf{Status:} #1\par}
\newcommand{\selector}[1]{\textbf{Selector:} #1\par}
\newcommand{\assuming}[1]{\textbf{Assuming:} #1\par}
\newcommand{\titlefield}[1]{\textbf{Title:} #1\par}
\let\titlemacro\title
\renewcommand{\title}[1]{\titlefield{#1}}
\renewenvironment{description}{\par\noindent\ignorespaces}{\par}
\makeatletter
\let\texlanguage\language
\def\languageenvname{language}
\def\language{\ifx\@currenvir\languageenvname\expandafter\@firstoftwo\else\expandafter\@secondoftwo\fi{\relax}{\texlanguage}}
\let\endlanguage\relax
\makeatother
\newenvironment{concept}{}{}
\newenvironment{succinctnessclaim}{}{}
\newenvironment{queryclaim}{}{}
\newenvironment{transformationclaim}{}{}
\newenvironment{batchclaim}{}{}
\titlemacro{Tractable Circuit Zoo: Transformation Support}
\date{\today}
\begin{document}
\maketitle

\begin{batchclaim}
\selector{\isin{\langref{cSDD}, \langref{d-SDNNF}, \langref{DNNF}, \langref{FBDD}, \langref{nFBDD}, \langref{nOBDD}, \langref{SDD}, \langref{SDNNF}, \langref{uFBDD}, \langref{uOBDD}, \langref{OBDD}, \langref{TDD}, \langref{dec-DNNF}, \langref{dec-SDNNF}}}
\operation{\ANDC}
\status{\nopolyunknownquasi}
\begin{description}
In the language \thislang, every clause has a polynomial-size representation. If \thislang supported $\land$C, then every \langref{CNF} could be compiled into \thislang by conjoining the translated clauses. This contradicts \nocompilespoly{\langref{CNF}}{\thislang} \citep{TCZ_Forthcoming}.

Alternatively, since \thislang does not support Bounded Conjunction in quasipolynomial time, intractability of Conjunction is immediate. \citep{Darwiche_2002}
\end{description}
\end{batchclaim}

\begin{batchclaim}
\selector{\isin{\langref{DNF}, \langfam{d-SDNNF}{T}, \langfam{SDD}{T}, \langfam{SDNNF}{T}, \langfam{dec-SDNNF}{T}, \langfam{cSDD}{T}, \langfam{TDD}{T}, \langfam{nOBDD}{<}, \langfam{uOBDD}{<}}}
\operation{\ANDC}
\status{\nopolyunknownquasi}
\begin{description}
In the language \thislang, every clause has a polynomial-size representation. If \thislang supported $\land$C, then every \langref{CNF} could be compiled into \thislang by conjoining the translated clauses. This contradicts \nocompilespoly{\langref{CNF}}{\thislang} \citep{TCZ_Forthcoming}. \citep{Darwiche_2002}
\end{description}
\end{batchclaim}

\begin{batchclaim}
\selector{\compilespoly{\langref{OBDD}}{\thislang}, \compilespoly{\thislang}{\langref{DNNF}}}
\operation{\ANDBC}
\status{\nopolyunknownquasi}
\begin{description}
There are Boolean functions $f,g:\{0,1\}^n\to\{0,1\}$, each computable by \langref{OBDD}{s} of size $O(n)$, such that any \langref{DNNF} computing $f\land g$ has size $2^{\Omega(\sqrt{n})}$ \citep{TCZ_Forthcoming}.

As a result, because \compilesquasi{\langref{OBDD}}{\thislang} and \compilesquasi{\thislang}{\langref{DNNF}}, \thislang cannot support bounded conjunction.
\end{description}
\end{batchclaim}

\begin{batchclaim}
\selector{\compilespoly{\langref{OBDD}}{\thislang}, \compilespoly{\thislang}{\langref{SDNNF}}}
\operation{\ORBC}
\status{\nopolyunknownquasi}
\begin{description}
There are Boolean functions $f,g:\{0,1\}^n\to\{0,1\}$, each computable by \langref{OBDD}{s} of size $O(n)$, such that any \langref{SDNNF} computing $f\lor g$ has size $2^{\Omega(\sqrt{n})}$ \citep{TCZ_Forthcoming}.

As a result, because \compilesquasi{\langref{OBDD}}{\thislang} and \compilesquasi{\thislang}{\langref{SDNNF}}, \thislang cannot support bounded disjunction in quasipolynomial time.
\end{description}
\end{batchclaim}

\begin{batchclaim}
\selector{\compilespoly{\langref{OBDD}}{\thislang}, \compilespoly{\thislang}{\langref{d-DNNF}}}
\operation{\ORBC}
\status{\nopolyunknownquasi}
\begin{description}
There are Boolean functions $f,g:\{0,1\}^n\to\{0,1\}$, each computable by \langref{OBDD}{s} of size $O(n)$, such that any \langref{d-DNNF} computing $f\lor g$ has size $2^{\Omega(\sqrt{n})}$ \citep{TCZ_Forthcoming}.

As a result, because \compilesquasi{\langref{OBDD}}{\thislang} and \compilesquasi{\thislang}{\langref{d-DNNF}}, \thislang cannot support bounded disjunction in quasipolynomial time.

Note that \citet[Theorem 9]{Bova_2016} proves the lower bound without explicitly noting the implication for Bounded Disjunction.
\end{description}
\end{batchclaim}

\begin{batchclaim}
\selector{\isin{\langfam{d-SDNNF}{T}, \langfam{uOBDD}{<}}}
\operation{\ORBC}
\status{\nopolyunknownquasi}
\begin{description}
\citet[Theorem 5]{Vinall-Smeeth_2024} gives unambiguous \langref{DNF} representations of functions $f$ and $g$ such that, measured in the size $s$ of the corresponding source representations, $f \lor g$ requires \langref{d-SDNNF} size $s^{\widetilde{\Omega}(\log s)}$. The unambiguous \langref{DNF} representations can be transformed into \thislang representations with only polynomial overhead, and every \thislang representation can be transformed into \langref{d-SDNNF} with only polynomial overhead. Thus polynomial-time bounded disjunction for \thislang would contradict the lower bound.
\end{description}
\end{batchclaim}

\begin{batchclaim}
\selector{\isin{\langref{cSDD}, \langref{d-DNNF}, \langref{dec-DNNF}, \langref{dec-SDNNF}, \langref{SDD}, \langref{uFBDD}, \langfam{d-SDNNF}{T}, \langfam{dec-SDNNF}{T}, \langfam{uOBDD}{<}}}
\operation{\SFO}
\status{\nopolyunknownquasi}
\begin{description}
If \thislang supported Singleton Forgetting in polynomial time, then given two \thislang circuits $f$ and $g$ we could form the decision $(x \land f) \lor (\neg x \land g)$ with only constant overhead and forget $x$, yielding $f \lor g$. This would give polynomial-time Bounded Disjunction for \thislang, contradicting the lower bound of \citet{Vinall-Smeeth_2024}. \citep{TCZ_Forthcoming}
\end{description}
\end{batchclaim}

\begin{batchclaim}
\selector{\compilespoly{\langref{DNF}}{\thislang}, \compilespoly{\thislang}{\langref{DNNF}}}
\operation{\NOTC}
\status{\nopolyunknownquasi}
\begin{description}
\citet[Lemma 5]{Bova_2014} gives 2-\langref{DNF} formulas $D$ whose negations require \langref{DNNF} size $2^{\Omega(\operatorname{size}(D))}$. If \thislang supported Negation in quasipolynomial time, then \compilesquasi{\langref{DNF}}{\thislang} and \compilesquasi{\thislang}{\langref{DNNF}} would yield quasipolynomial-size \langref{DNNF} formulas for those negations, contradicting the lower bound.
\end{description}
\end{batchclaim}

\begin{batchclaim}
\selector{\isin{\langref{cSDD}, \langref{d-SDNNF}, \langref{FBDD}, \langref{uOBDD}, \langref{OBDD}, \langref{TDD}, \langref{dec-DNNF}, \langref{dec-SDNNF}}}
\operation{\ORC}
\status{\nopolyunknownquasi}
\begin{description}
In the language \thislang, every term has a polynomial-size representation. If \thislang supported $\lor$C, then every \langref{DNF} could be compiled into \thislang by disjoining the translated terms. This contradicts \nocompilespoly{\langref{DNF}}{\thislang} \citep{TCZ_Forthcoming}.

Alternatively, since \thislang does not support Bounded Disjunction in quasipolynomial time, intractability of Disjunction is immediate. \citep{Darwiche_2002}
\end{description}
\end{batchclaim}

\begin{batchclaim}
\selector{\isin{\langfam{d-SDNNF}{T}, \langfam{OBDD}{<}, \langfam{SDD}{T}, \langfam{dec-SDNNF}{T}, \langfam{cSDD}{T}, \langfam{TDD}{T}, \langfam{uOBDD}{<}}}
\operation{\ORC}
\status{\nopolyunknownquasi}
\begin{description}
In the language \thislang, every term has a polynomial-size representation. If \thislang supported $\lor$C, then every \langref{DNF} could be compiled into \thislang by disjoining the translated terms. This contradicts \nocompilespoly{\langref{DNF}}{\thislang} \citep{TCZ_Forthcoming}. \citep{Darwiche_2002}
\end{description}
\end{batchclaim}

\begin{batchclaim}
\selector{\isin{\langref{CNF}, \langref{NNF}}}
\operation{\FO}
\status{\nopolyunknownquasi}
\assuming{P \neq NP}
\begin{description}
If \thislang supported FO, then the consistency of any \langref{CNF} representation $\Sigma$ could be tested in polynomial time: represent $\Sigma$ in \thislang, compute $\exists\text{Vars}(\Sigma).\Sigma$, and check the resulting constant representation. This would decide \langref{CNF} consistency in polynomial time, implying P = NP \citep{Darwiche_2002}.
\end{description}
\end{batchclaim}

\begin{batchclaim}
\selector{\isin{\langref{OBDD}, \langref{TDD}, \langfam{OBDD}{<}, \langfam{SDD}{T}, \langfam{cSDD}{T}, \langfam{TDD}{T}}}
\operation{\FO}
\status{\nopolyunknownquasi}
\begin{description}
A quasipolynomial-time forgetting procedure for \thislang would turn every \langref{DNF} representation into an equivalent \thislang representation of quasipolynomial size in the input \langref{DNF} after introducing fresh selectors and forgetting them. This contradicts \nocompilespoly{\langref{DNF}}{\thislang} \citep{TCZ_Forthcoming}. \citep{Darwiche_2002,VanDenBroeck_2015}
\end{description}
\end{batchclaim}

\begin{transformationclaim}
\claimlanguage{\langref{CNF}}
\operation{\ANDC}
\status{\poly}
\begin{description}
Conjunction of \langref{CNF} representation is trivially another \langref{CNF} representation: simply take the union of all clauses \citep{Darwiche_2002}.
\end{description}
\end{transformationclaim}

\begin{transformationclaim}
\claimlanguage{\langref{CNF}}
\operation{\ORBC}
\status{\poly}
\begin{description}
The disjunction of two \langref{CNF} representations $\alpha_1 \lor \alpha_2$ can be computed by taking all cross products of one clause from $\alpha_1$ and one from $\alpha_2$, removing redundant literals and valid clauses. The result is a \langref{CNF} of polynomial size \citep{Darwiche_2002}.
\end{description}
\end{transformationclaim}

\begin{transformationclaim}
\claimlanguage{\langref{CNF}}
\operation{\ORC}
\status{\nopolyunknownquasi}
\begin{description}
It is straightforward to convert any term into an equivalent \langref{CNF} formula in polynomial time. Since specific \langref{DNF} formulas cannot be turned into equivalent \langref{CNF} formulas in polynomial space (\langref{DNF} $\not\leq$ \langref{CNF}), \langref{CNF} does not satisfy $\lor$C \citep{Darwiche_2002}.
\end{description}
\end{transformationclaim}

\begin{transformationclaim}
\claimlanguage{\langref{CNF}}
\operation{\CD}
\status{\poly}
\begin{description}
CD is trivially satisfied: replacing literals of the conditioning term by Boolean constants in a \langref{CNF} and removing valid clauses preserves the \langref{CNF} property \citep{Darwiche_2002}.
\end{description}
\end{transformationclaim}

\begin{transformationclaim}
\claimlanguage{\langref{cSDD}}
\operation{\NOTC}
\status{\poly}
\begin{description}
Every \langref{cSDD} respects some vtree $T$. Apply the fixed-vtree negation procedure for \langfam{cSDD}{T}: \citet{VanDenBroeck_2015} proved that reduced \langfam{SDD}{T}{s} support negation in polynomial time by passing negation through the subs. The result is still a \langfam{cSDD}{T} formula, hence a \langref{cSDD} formula.
\end{description}
\end{transformationclaim}

\begin{transformationclaim}
\claimlanguage{\langref{d-DNNF}}
\operation{\CD}
\status{\poly}
\begin{description}
Conditioning preserves both decomposability and determinism. If $\alpha \land \beta \models \bot$ then $(\alpha|\gamma) \land (\beta|\gamma) \models \bot$ for any consistent term $\gamma$, since $((\alpha \land \beta)|\gamma) \land \gamma \models \bot$ and $\gamma$ is consistent \citep{Darwiche_2002}.
\end{description}
\end{transformationclaim}

\begin{transformationclaim}
\claimlanguage{\langref{d-SDNNF}}
\operation{\CD}
\status{\poly}
\begin{description}
Conditioning preserves both structured decomposability and determinism. The vtree structure is maintained since $\text{Vars}(\Sigma|\gamma) \subseteq \text{Vars}(\Sigma)$, and mutual exclusivity of $\lor$-children is preserved under conditioning \citet[Proposition 4]{Pipatsrisawat_2008}.
\end{description}
\end{transformationclaim}

\begin{transformationclaim}
\claimlanguage{\langref{d-SDNNF}}
\operation{\FO}
\status{\nopolyunknownquasi}
\begin{description}
If \langref{d-SDNNF} satisfied FO, then every \langref{DNF} could be compiled to an equivalent \langref{d-SDNNF} of size polynomial in the input \langref{DNF} (using fresh selector variables and applying FO). Since specific \langref{DNF} formulas require exponential-size \langref{d-SDNNF} representations (unconditionally), this is impossible. This unconditional result applies to the fixed-vtree case \citet[Proposition 4]{Pipatsrisawat_2008}.
\end{description}
\end{transformationclaim}

\begin{transformationclaim}
\claimlanguage{\langref{d-SDNNF}}
\operation{\NOTC}
\status{\nopolyunknownquasi}
\begin{description}
\langref{d-SDNNF} is not closed under negation in polytime. \citet{Vinall-Smeeth_2024} proves that for every $n \in \mathbb{N}$, there exists a Boolean function $f$ with an equivalent \langref{d-SDNNF} of size $n$ such that any \langref{SDNNF} equivalent to $\neg f$ has size $n^{\Omega(\tilde{\log}\, n)}$. This super-polynomial lower bound holds for every vtree (not just a single tree), so the result applies to \langref{d-SDNNF} (not just \langfam{d-SDNNF}{T}). Whether quasipolynomial closure holds remains open.
\end{description}
\end{transformationclaim}

\begin{transformationclaim}
\claimlanguage{\langref{d-SDNNF}}
\operation{\SFO}
\status{\nopolyunknownquasi}
\begin{description}
\citet[Corollary 1]{Vinall-Smeeth_2024} gives \langref{d-SDNNF} formulas of size $s$ whose existential quantification of one variable requires \langref{d-SDNNF} size $s^{\widetilde{\Omega}(\log s)}$. Hence \langref{d-SDNNF} does not support singleton forgetting in polynomial time.
\end{description}
\end{transformationclaim}

\begin{transformationclaim}
\claimlanguage{\langfam{d-SDNNF}{T}}
\operation{\ANDBC}
\status{\poly}
\begin{description}
Bounded conjunction of two \langfam{d-SDNNF}{T} formulas can be performed in polytime via the Apply algorithm, preserving both determinism and structured decomposability \citet[Proposition 4]{Pipatsrisawat_2008}.
\end{description}
\end{transformationclaim}

\begin{transformationclaim}
\claimlanguage{\langfam{d-SDNNF}{T}}
\operation{\CD}
\status{\poly}
\begin{description}
Conditioning preserves both decomposability and determinism, and it preserves the fixed vtree because restricting by a term only removes variables or replaces literals by constants. Hence the conditioned formula remains in \langfam{d-SDNNF}{T} \citet[Proposition 4]{Pipatsrisawat_2008}.
\end{description}
\end{transformationclaim}

\begin{transformationclaim}
\claimlanguage{\langref{DNF}}
\operation{\ANDBC}
\status{\poly}
\begin{description}
The conjunction of two \langref{DNF} representations $\alpha_1 \land \alpha_2$ can be computed by taking all cross products of one term from $\alpha_1$ and one from $\alpha_2$, removing redundant literals and inconsistent terms. The result is a \langref{DNF} of polynomial size \citep{Darwiche_2002}.
\end{description}
\end{transformationclaim}

\begin{transformationclaim}
\claimlanguage{\langref{DNF}}
\operation{\ORC}
\status{\poly}
\begin{description}
Disjunction of \langref{DNF} representations is trivially another \langref{DNF} representation: simply take the union of all terms \citep{Darwiche_2002}.
\end{description}
\end{transformationclaim}

\begin{transformationclaim}
\claimlanguage{\langref{DNF}}
\operation{\FO}
\status{\poly}
\begin{description}
\langref{DNF} satisfies FO \citep{Lang_2003}. Forgetting variables from a \langref{DNF} representation can be done in polytime \citep{Darwiche_2002}.
\end{description}
\end{transformationclaim}

\begin{transformationclaim}
\claimlanguage{\langref{DNNF}}
\operation{\ORC}
\status{\poly}
\begin{description}
Disjunction trivially preserves decomposability since decomposability only constrains $\land$-nodes. \citep{Darwiche_2002}.
\end{description}
\end{transformationclaim}

\begin{transformationclaim}
\claimlanguage{\langref{DNNF}}
\operation{\CD}
\status{\poly}
\begin{description}
Conditioning preserves decomposability: if $\alpha$ and $\beta$ do not share variables, then $\alpha|\gamma$ and $\beta|\gamma$ do not share variables either, since $\text{Vars}(\alpha|\gamma) \subseteq \text{Vars}(\alpha)$ \citep{Darwiche_2002}.
\end{description}
\end{transformationclaim}

\begin{transformationclaim}
\claimlanguage{\langref{DNNF}}
\operation{\FO}
\status{\poly}
\begin{description}
\langref{DNNF} satisfies FO: forgetting a set of variables in a \langref{DNNF} representation can be performed in polynomial time while preserving decomposability \citep{Darwiche_2001a} \citep{Darwiche_2002}.
\end{description}
\end{transformationclaim}

\begin{transformationclaim}
\claimlanguage{\langref{FBDD}}
\operation{\NOTC}
\status{\poly}
\begin{description}
Negation of an \langref{FBDD} is achieved by switching the labels of the 0-sink and 1-sink. The decision property and read-once property are preserved \citep{Darwiche_2002}.
\end{description}
\end{transformationclaim}

\begin{transformationclaim}
\claimlanguage{\langref{FBDD}}
\operation{\CD}
\status{\poly}
\begin{description}
Conditioning is easy on branching programs: replace each decision node labeled by a conditioned variable with the child consistent with the conditioning term. This preserves the free read-once property \citep{Darwiche_2002}.
\end{description}
\end{transformationclaim}

\begin{transformationclaim}
\claimlanguage{\langref{FBDD}}
\operation{\SFO}
\status{\nopolyunknownquasi}
\begin{description}
Let $ROW_n$ and $COL_n$ test whether an $n \times n$ Boolean matrix has an all-1 row or an all-1 column. The selector construction $(s \land ROW_n) \lor (\neg s \land COL_n)$ has polynomial-size \langref{FBDD}{s}: the top variable $s$ chooses a rowwise or columnwise computation. Forgetting $s$ yields $ROW_n \lor COL_n$, whose \langref{FBDD} size is $\Omega(n^{-7/2}2^n)$ by \citet[Theorem 6.2.13]{Wegener_2000}.
\end{description}
\end{transformationclaim}

\begin{transformationclaim}
\claimlanguage{\langref{IP}}
\operation{\NOTC}
\status{\nopolyunknownquasi}
\begin{description}
Dual of \langref{PI}: $\Sigma_n = \bigvee_{i=0}^{n-1}(x_{2i} \land x_{2i+1})$ is in \langref{IP} form but has $2^n$ prime implicates. Its negation has $2^n$ prime implicants. Hence \langref{IP} cannot satisfy $\neg$C \citep{Darwiche_2002}.
\end{description}
\end{transformationclaim}

\begin{transformationclaim}
\claimlanguage{\langref{IP}}
\operation{\ANDBC}
\status{\poly}
\begin{description}
$\text{IP}(\alpha_1 \land \alpha_2) = \max(\{\beta_1 \land \beta_2 \mid \beta_1 \in \text{IP}(\alpha_1), \beta_2 \in \text{IP}(\alpha_2)\}, \models)$ (dual of Proposition 40, \citet{Marquis_2000}). The maximization removes subsuming terms, yielding an \langref{IP} form in polytime \citep{Darwiche_2002,Marquis_2000}.
\end{description}
\end{transformationclaim}

\begin{transformationclaim}
\claimlanguage{\langref{IP}}
\operation{\ANDC}
\status{\nopolyunknownquasi}
\begin{description}
It is straightforward to convert a clause into an equivalent \langref{IP} formula in polynomial time. Since specific \langref{CNF} formulas cannot be turned into equivalent \langref{IP} formulas in polynomial space (\langref{CNF} $\not\leq$ \langref{IP}), \langref{IP} does not satisfy $\land$C \citep{Darwiche_2002}.
\end{description}
\end{transformationclaim}

\begin{transformationclaim}
\claimlanguage{\langref{IP}}
\operation{\ORBC}
\status{\nopolyunknownquasi}
\begin{description}
Combining prime implicants via disjunction requires distributing and re-minimizing terms, which can lead to an exponential explosion of new prime implicants. Let $\alpha_1 = \bigwedge_{i=1}^k p_i$ (one prime implicant) and $\alpha_2 = \bigvee_{i=1}^k \bigvee_{j=1}^m (\neg p_i \land q_{i,j})$ ($mk$ prime implicants). Then $\alpha_1 \lor \alpha_2$ has $(m+1)^k + mk$ prime implicants, an exponential blowup \citep{Chandra_1978} \citep{Darwiche_2002}.
\end{description}
\end{transformationclaim}

\begin{transformationclaim}
\claimlanguage{\langref{IP}}
\operation{\CD}
\status{\poly}
\begin{description}
Given an \langref{IP} formula $\Sigma = \bigvee \gamma_i$, conditioning on $\gamma$ gives $\bigvee (\gamma_i|\gamma)$. Keeping only the logically weakest surviving terms yields a prime implicants formula equivalent to $\Sigma|\gamma$. This requires $O(n^2)$ entailment tests among terms, each polytime \citep{Darwiche_2002}.
\end{description}
\end{transformationclaim}

\begin{transformationclaim}
\claimlanguage{\langref{IP}}
\operation{\SFO}
\status{\nopolyunknownquasi}
\begin{description}
The number of prime implicants of $\exists x.\Sigma$ can be exponentially greater than the number of prime implicants of $\Sigma$ \citep{Chandra_1978}. A specific construction with $mk+1$ prime implicants produces $(m+1)^k + mk$ prime implicants after forgetting one variable \citep{Darwiche_2002}.
\end{description}
\end{transformationclaim}

\begin{transformationclaim}
\claimlanguage{\langref{MODS}}
\operation{\NOTC}
\status{\nopolyunknownquasi}
\begin{description}
$\Sigma = \bigwedge_{i=1}^n x_i$ has only one model over $\{x_1,\ldots,x_n\}$ but its negation $\neg\Sigma$ has $2^n - 1$ models. Hence \langref{MODS} cannot satisfy $\neg$C \citep{Darwiche_2002}.
\end{description}
\end{transformationclaim}

\begin{transformationclaim}
\claimlanguage{\langref{MODS}}
\operation{\ANDBC}
\status{\poly}
\begin{description}
The bounded conjunction of two \langref{MODS} formulas (over the same variables) is simply taking the intersection of their model sets. If the variables differ, it is the intersection of their cylindrical extensions. The result is a \langref{MODS} representation in polytime \citep{Darwiche_2002}.
\end{description}
\end{transformationclaim}

\begin{transformationclaim}
\claimlanguage{\langref{MODS}}
\operation{\ANDC}
\status{\nopolyunknownquasi}
\begin{description}
Let $\Sigma_i = (x_{i,1} \lor x_{i,2})$ for $i = 1,\ldots,n$. Each $\Sigma_i$ has 3 models. But $\Sigma = \bigwedge_{i=1}^n \Sigma_i$ has $3^n$ models, so no polynomial-size \langref{MODS} representation exists \citep{Darwiche_2002}.
\end{description}
\end{transformationclaim}

\begin{transformationclaim}
\claimlanguage{\langref{MODS}}
\operation{\ORBC}
\status{\nopolyunknownquasi}
\begin{description}
Let $\alpha_1 = \bigwedge_{i=1}^n x_i$ (1 model) and $\alpha_2 = y$ (1 model). Then $\alpha_1 \lor \alpha_2$ has $2^n + 1$ models over $\text{Vars}(\alpha_1) \cup \text{Vars}(\alpha_2)$. Since this already uses the disjunction of just two \langref{MODS} representations, \langref{MODS} does not satisfy $\lor$BC \citep{Darwiche_2002}.
\end{description}
\end{transformationclaim}

\begin{transformationclaim}
\claimlanguage{\langref{MODS}}
\operation{\FO}
\status{\poly}
\begin{description}
Given a \langref{MODS} formula $\Sigma$ and a set $X$, the \langref{MODS} representation of $\exists X.\Sigma$ is obtained by removing every leaf node labeled by a literal $x$ or $\neg x$ with $x \in X$ (Propositions 18 and 20, Lang et al. 2000). See also the polytime FO operation on \langref{DNNF} from \citep{Darwiche_2001a} \citep{Darwiche_2002}.
\end{description}
\end{transformationclaim}

\begin{transformationclaim}
\claimlanguage{\langref{nFBDD}}
\operation{\ORC}
\status{\poly}
\begin{description}
Disjunction of \langref{nFBDD} formulas is obtained by adding a top-level nondeterministic $\lor$-node whose children are the input formulas. This preserves the free read-once property. \citep{TCZ_Forthcoming}
\end{description}
\end{transformationclaim}

\begin{transformationclaim}
\claimlanguage{\langref{nFBDD}}
\operation{\CD}
\status{\poly}
\begin{description}
Conditioning is easy on branching programs: replace each decision node labeled by a conditioned variable with the child consistent with the conditioning term. This preserves the free read-once property. \citet{Wegener_2000} gives this argument for EXOR-OBDDs in Theorem 10.5.1.
\end{description}
\end{transformationclaim}

\begin{transformationclaim}
\claimlanguage{\langref{nFBDD}}
\operation{\FO}
\status{\poly}
\begin{description}
To forget a variable, replace each decision node labeled by that variable with a nondeterministic $\lor$-node over its two children. Repeating this for all forgotten variables preserves the free read-once property and computes existential quantification. \citep{TCZ_Forthcoming}
\end{description}
\end{transformationclaim}

\begin{transformationclaim}
\claimlanguage{\langref{NNF}}
\operation{\NOTC}
\status{\poly}
\begin{description}
\langref{NNF} trivially satisfies $\neg$C since any propositional representation can be converted to \langref{NNF} by pushing negations to literals and this preserves the \langref{NNF} property \citep{Darwiche_2002}.
\end{description}
\end{transformationclaim}

\begin{transformationclaim}
\claimlanguage{\langref{NNF}}
\operation{\ANDC}
\status{\poly}
\begin{description}
Taking the conjunction of \langref{NNF} representations $\alpha_1 \land \alpha_2 \land \ldots$ trivially produces another \langref{NNF} representation, since \langref{NNF} places no restrictions on $\land$-nodes \citep{Darwiche_2002}.
\end{description}
\end{transformationclaim}

\begin{transformationclaim}
\claimlanguage{\langref{NNF}}
\operation{\CD}
\status{\poly}
\begin{description}
CD is trivially satisfied by \langref{NNF}: replacing literals of the conditioning term by Boolean constants preserves the \langref{NNF} property \citep{Darwiche_2002}.
\end{description}
\end{transformationclaim}

\begin{transformationclaim}
\claimlanguage{\langref{nOBDD}}
\operation{\CD}
\status{\poly}
\begin{description}
Conditioning is easy on branching programs: replace each decision node labeled by a conditioned variable with the child consistent with the conditioning term. This preserves the ordered decision-diagram property. \citet{Wegener_2000} gives this argument for EXOR-OBDDs in Theorem 10.5.1.
\end{description}
\end{transformationclaim}

\begin{transformationclaim}
\claimlanguage{\langref{nOBDD}}
\operation{\FO}
\status{\poly}
\begin{description}
To forget a variable, replace each decision node labeled by that variable with a nondeterministic $\lor$-node over its two children. Repeating this for all forgotten variables preserves the ordered decision-diagram property and computes existential quantification. \citep{TCZ_Forthcoming}
\end{description}
\end{transformationclaim}

\begin{transformationclaim}
\claimlanguage{\langfam{OBDD}{<}}
\operation{\NOTC}
\status{\poly}
\begin{description}
Negation of an \langfam{OBDD}{<} is achieved by switching the labels of the 0-sink and 1-sink. The result is an \langfam{OBDD}{<} of identical size \citep{Bryant_1986}.
\end{description}
\end{transformationclaim}

\begin{transformationclaim}
\claimlanguage{\langfam{OBDD}{<}}
\operation{\ANDBC}
\status{\poly}
\begin{description}
The Apply algorithm computes the conjunction of two \langfam{OBDD}{<} representations sharing the same variable ordering in polynomial time. The result is a reduced \langfam{OBDD}{<} \citep{Bryant_1986}.
\end{description}
\end{transformationclaim}

\begin{transformationclaim}
\claimlanguage{\langfam{OBDD}{<}}
\operation{\ORBC}
\status{\poly}
\begin{description}
The Apply algorithm computes the disjunction of two \langfam{OBDD}{<} representations sharing the same variable ordering in polynomial time. The result is a reduced \langfam{OBDD}{<} \citep{Bryant_1986}.
\end{description}
\end{transformationclaim}

\begin{transformationclaim}
\claimlanguage{\langfam{OBDD}{<}}
\operation{\CD}
\status{\poly}
\begin{description}
Conditioning is easy on \langfam{OBDD}{<}: replace each decision node labeled by a conditioned variable with the child consistent with the conditioning term. The fixed order $<$ is preserved \citep{Bryant_1986}.
\end{description}
\end{transformationclaim}

\begin{transformationclaim}
\claimlanguage{\langref{PI}}
\operation{\NOTC}
\status{\nopolyunknownquasi}
\begin{description}
The formula $\Sigma_n = \bigwedge_{i=0}^{n-1}(x_{2i} \lor x_{2i+1})$ is in \langref{PI} form but has $2^n$ prime implicants. Its negation has $2^n$ prime implicates. Hence \langref{PI} cannot satisfy $\neg$C \citep{Darwiche_2002}.
\end{description}
\end{transformationclaim}

\begin{transformationclaim}
\claimlanguage{\langref{PI}}
\operation{\ANDBC}
\status{\nopolyunknownquasi}
\begin{description}
Combining prime implicates via conjunction requires distributing and re-minimizing clauses, which can lead to an exponential explosion of new prime implicates. Let $\alpha_1 = \bigvee_{i=1}^k p_i$ (one prime implicate) and $\alpha_2 = \bigwedge_{i=1}^k \bigwedge_{j=1}^m (\neg p_i \lor q_{i,j})$ ($mk$ prime implicates). Then $\alpha_1 \land \alpha_2$ has $(m+1)^k + mk$ prime implicates, an exponential blowup \citep{Chandra_1978} \citep{Darwiche_2002}.
\end{description}
\end{transformationclaim}

\begin{transformationclaim}
\claimlanguage{\langref{PI}}
\operation{\ORBC}
\status{\poly}
\begin{description}
$\text{PI}(\alpha_1 \lor \alpha_2) = \min(\{\beta_1 \lor \beta_2 \mid \beta_1 \in \text{PI}(\alpha_1), \beta_2 \in \text{PI}(\alpha_2)\}, \models)$ (Proposition 40, \citet{Marquis_2000}). The minimization removes subsumed clauses, yielding a \langref{PI} form in polytime \citep{Darwiche_2002,Marquis_2000}.
\end{description}
\end{transformationclaim}

\begin{transformationclaim}
\claimlanguage{\langref{PI}}
\operation{\ORC}
\status{\nopolyunknownquasi}
\begin{description}
It is straightforward to convert any term into an equivalent \langref{PI} formula in polynomial time. Since specific \langref{DNF} formulas cannot be turned into equivalent \langref{PI} formulas in polynomial space (\langref{DNF} $\not\leq$ \langref{PI}), \langref{PI} does not satisfy $\lor$C \citep{Darwiche_2002}.
\end{description}
\end{transformationclaim}

\begin{transformationclaim}
\claimlanguage{\langref{PI}}
\operation{\CD}
\status{\poly}
\begin{description}
The prime implicates of $\Sigma \land \gamma$ can be computed in polytime when $\Sigma$ is in \langref{PI} form and $\gamma$ is a term (Proposition 36 and the following paragraph, \citet{Marquis_2000}). Since \langref{PI} satisfies FO, the prime implicates of $\Sigma|\gamma \equiv \exists\text{Vars}(\gamma).(\Sigma \land \gamma)$ follow in polytime \citep{Darwiche_2002,Marquis_2000}.
\end{description}
\end{transformationclaim}

\begin{transformationclaim}
\claimlanguage{\langref{PI}}
\operation{\FO}
\status{\poly}
\begin{description}
The prime implicates of $\exists X.\Sigma$ are exactly the prime implicates of $\Sigma$ that do not contain any variable from $X$ (Proposition 55, \citet{Marquis_2000}). Hence FO for \langref{PI} simply filters out clauses containing forgotten variables \citep{Darwiche_2002,Marquis_2000}.
\end{description}
\end{transformationclaim}

\begin{transformationclaim}
\claimlanguage{\langref{SDD}}
\operation{\NOTC}
\status{\poly}
\begin{description}
An \langref{SDD} can be negated in polytime by applying exclusive-or with $\top$ using the Apply algorithm \citep{Darwiche_2011}.
\end{description}
\end{transformationclaim}

\begin{transformationclaim}
\claimlanguage{\langref{SDD}}
\operation{\CD}
\status{\poly}
\begin{description}
Conditioning an \langref{SDD} $f$ on a literal $l$ is equivalent to computing $f \land l$. Since any literal is a trivially small \langref{SDD} and \langref{SDD} supports $\land$BC in polytime via Apply, conditioning is polytime \citep{Darwiche_2011}.
\end{description}
\end{transformationclaim}

\begin{transformationclaim}
\claimlanguage{\langfam{SDD}{T}}
\operation{\NOTC}
\status{\poly}
\begin{description}
Negation of an \langfam{SDD}{T} can be computed in polytime by applying exclusive-or with $\top$ using Apply. Since $\top$ trivially respects any fixed vtree $T$, the result remains an \langfam{SDD}{T} over the same vtree \citep{Darwiche_2011}.
\end{description}
\end{transformationclaim}

\begin{transformationclaim}
\claimlanguage{\langfam{SDD}{T}}
\operation{\ANDBC}
\status{\poly}
\begin{description}
The Apply algorithm combines two \langfam{SDD}{T} representations using any Boolean operator in polytime. In particular, two \langfam{SDD}{T}{s} can be conjoined in polytime via Apply \citep{Darwiche_2011}.
\end{description}
\end{transformationclaim}

\begin{transformationclaim}
\claimlanguage{\langfam{SDD}{T}}
\operation{\ORBC}
\status{\poly}
\begin{description}
The Apply algorithm combines two \langfam{SDD}{T} representations using any Boolean operator in polytime. In particular, two \langfam{SDD}{T}{s} can be disjoined in polytime via Apply \citep{Darwiche_2011}.
\end{description}
\end{transformationclaim}

\begin{transformationclaim}
\claimlanguage{\langfam{SDD}{T}}
\operation{\CD}
\status{\poly}
\begin{description}
Conditioning an \langfam{SDD}{T} representation $f$ on a literal $l$ is equivalent to computing $f \land l$. Since a literal is a trivially small \langfam{SDD}{T} respecting the same fixed vtree, and \langfam{SDD}{T} supports $\land$BC in polytime via Apply, conditioning is polynomial-time as well \citep{Darwiche_2011}.
\end{description}
\end{transformationclaim}

\begin{transformationclaim}
\claimlanguage{\langref{SDNNF}}
\operation{\CD}
\status{\poly}
\begin{description}
Conditioning preserves structured decomposability: replacing conditioned literals by constants and simplifying can only remove variables from subformulas, so every decomposable split continues to respect the original vtree. Therefore the result remains an \langref{SDNNF} formula of size polynomial in the input formula \citet[Proposition 4]{Pipatsrisawat_2008}.
\end{description}
\end{transformationclaim}

\begin{transformationclaim}
\claimlanguage{\langref{SDNNF}}
\operation{\FO}
\status{\poly}
\begin{description}
Forgetting preserves structured decomposability. Replacing all occurrences of literal $x$ and $\neg x$ with $\top$ in an \langref{SDNNF} preserves the vtree structure, since the decomposition at $\land$-nodes remains valid after variable removal. Inherited from \langref{DNNF} \citet[Proposition 4]{Pipatsrisawat_2008}.
\end{description}
\end{transformationclaim}

\begin{transformationclaim}
\claimlanguage{\langfam{SDNNF}{T}}
\operation{\ANDBC}
\status{\poly}
\begin{description}
Bounded conjunction of two \langfam{SDNNF}{T} formulas can be performed in polytime via the Apply algorithm. The fixed vtree governs variable partitioning at $\land$-nodes, enabling an efficient recursive traversal \citet[Proposition 4]{Pipatsrisawat_2008}.
\end{description}
\end{transformationclaim}

\begin{transformationclaim}
\claimlanguage{\langfam{SDNNF}{T}}
\operation{\ORC}
\status{\poly}
\begin{description}
Disjunction of \langfam{SDNNF}{T} formulas preserves structured decomposability, since decomposability only constrains $\land$-nodes and a top-level $\lor$-node introduces no new $\land$-node constraints \citet[Proposition 4]{Pipatsrisawat_2008}.
\end{description}
\end{transformationclaim}

\begin{transformationclaim}
\claimlanguage{\langfam{SDNNF}{T}}
\operation{\CD}
\status{\poly}
\begin{description}
Conditioning preserves decomposability with respect to the same fixed vtree: after replacing conditioned literals by constants, each remaining subformula mentions a subset of its previous variables, so every $\land$-node still decomposes according to the original vtree split \citet[Proposition 4]{Pipatsrisawat_2008}.
\end{description}
\end{transformationclaim}

\begin{transformationclaim}
\claimlanguage{\langfam{SDNNF}{T}}
\operation{\FO}
\status{\poly}
\begin{description}
Darwiche's \langref{DNNF} projection operation for forgetting variables preserves decomposability with respect to any fixed vtree. Applying that operation to a \langfam{SDNNF}{T} formula therefore yields an equivalent \langfam{SDNNF}{T} formula in polynomial time \citet[Proposition 4]{Pipatsrisawat_2008}.
\end{description}
\end{transformationclaim}

\begin{transformationclaim}
\claimlanguage{\langref{uFBDD}}
\operation{\CD}
\status{\poly}
\begin{description}
Conditioning is easy on branching programs: replace each decision node labeled by a conditioned variable with the child consistent with the conditioning term. This preserves the free read-once property and cannot create two accepting paths for a remaining assignment. \citet{Wegener_2000} gives this argument for EXOR-OBDDs in Theorem 10.5.1.
\end{description}
\end{transformationclaim}

\begin{transformationclaim}
\claimlanguage{\langref{uOBDD}}
\operation{\CD}
\status{\poly}
\begin{description}
Conditioning is easy on branching programs: replace each decision node labeled by a conditioned variable with the child consistent with the conditioning term. This preserves the ordered decision-diagram property and cannot create two accepting paths for a remaining assignment. \citet{Wegener_2000} gives this argument for EXOR-OBDDs in Theorem 10.5.1.
\end{description}
\end{transformationclaim}

\begin{transformationclaim}
\claimlanguage{\langref{uOBDD}}
\operation{\NOTC}
\status{\nopolyunknownquasi}
\begin{description}
\citet[Theorem 2]{Vinall-Smeeth_2024} gives unambiguous \langref{DNF} representations whose negations require quasipolynomially larger \langref{d-SDNNF} size, measured against the source representation size. An unambiguous \langref{DNF} can be transformed into a \langref{uOBDD} with only polynomial overhead, and every \langref{uOBDD} can be transformed into \langref{d-SDNNF} with only polynomial overhead. Thus polynomial-time negation for \langref{uOBDD} would contradict the lower bound.
\end{description}
\end{transformationclaim}

\begin{transformationclaim}
\claimlanguage{\langref{uOBDD}}
\operation{\SFO}
\status{\nopolyunknownquasi}
\begin{description}
The proof of \citet[Corollary 1]{Vinall-Smeeth_2024} shows the intractability of SFO for \langref{d-SDNNF}, but also immediately applies to \langref{uOBDD}.
\end{description}
\end{transformationclaim}

\begin{transformationclaim}
\claimlanguage{\langref{OBDD}}
\operation{\NOTC}
\status{\poly}
\begin{description}
Negation of an \langref{OBDD} is achieved by switching the labels of the 0-sink and 1-sink. The ordered variable property is preserved since the single ordering is not affected \citep{Darwiche_2002}.
\end{description}
\end{transformationclaim}

\begin{transformationclaim}
\claimlanguage{\langref{OBDD}}
\operation{\CD}
\status{\poly}
\begin{description}
Since \langref{OBDD} $\subseteq$ \langref{FBDD} and conditioning preserves the variable ordering, \langref{OBDD} also satisfies CD \citep{Darwiche_2002}.
\end{description}
\end{transformationclaim}

\begin{transformationclaim}
\claimlanguage{\langref{OBDD}}
\operation{\SFO}
\status{\poly}
\begin{description}
Since only one \langref{OBDD} representation is considered in the transformation, and \langfam{OBDD}{<} satisfies SFO, \langref{OBDD} also satisfies SFO \citep{Darwiche_2002}.
\end{description}
\end{transformationclaim}

\begin{transformationclaim}
\claimlanguage{\langfam{TDD}{T}}
\operation{\CD}
\status{\poly}
\begin{description}
Conditioning a \langfam{TDD}{T} on a literal can be performed in polynomial time while preserving the same vtree, as stated in \citet[Theorem 4]{Capelli_2026}.
\end{description}
\end{transformationclaim}

\begin{transformationclaim}
\claimlanguage{\langfam{TDD}{T}}
\operation{\ANDBC}
\status{\poly}
\begin{description}
Bounded conjunction is polynomial-time for \langfam{TDD}{T}{s} respecting the same vtree by the apply construction of \citet[Theorem 4]{Capelli_2026}.
\end{description}
\end{transformationclaim}

\begin{transformationclaim}
\claimlanguage{\langfam{TDD}{T}}
\operation{\NOTC}
\status{\poly}
\begin{description}
\citet[Theorem 4]{Capelli_2026} proves that any \langfam{TDD}{T} can be negated in polynomial time while preserving the same vtree.
\end{description}
\end{transformationclaim}

\begin{transformationclaim}
\claimlanguage{\langref{TDD}}
\operation{\CD}
\status{\poly}
\begin{description}
\citet[Theorem 4]{Capelli_2026} proves that any \langref{TDD} can be conditioned in polynomial time.
\end{description}
\end{transformationclaim}

\begin{transformationclaim}
\claimlanguage{\langref{TDD}}
\operation{\SFO}
\status{\poly}
\begin{description}
\citet[Theorem 4]{Capelli_2026} proves that any \langref{TDD} can be existentially quantified on a single variable in polynomial time.
\end{description}
\end{transformationclaim}

\begin{transformationclaim}
\claimlanguage{\langref{TDD}}
\operation{\NOTC}
\status{\poly}
\begin{description}
\citet[Theorem 4]{Capelli_2026} proves that any \langref{TDD} can be negated in polynomial time.
\end{description}
\end{transformationclaim}

\begin{transformationclaim}
\claimlanguage{\langref{dec-DNNF}}
\operation{\CD}
\status{\poly}
\begin{description}
A \langref{dec-DNNF} can be viewed as a decision diagram with additional decomposable $\land$-nodes. Conditioning is easy on branching programs: replace each decision node labeled by a conditioned variable with the child consistent with the conditioning term. The extra $\land$-nodes cause no problem: recurse through them, and decomposability is preserved because conditioning only removes variables from subcircuits. \citet{Wegener_2000} gives this argument for EXOR-OBDDs in Theorem 10.5.1.
\end{description}
\end{transformationclaim}

\begin{transformationclaim}
\claimlanguage{\langref{dec-SDNNF}}
\operation{\CD}
\status{\poly}
\begin{description}
A \langref{dec-SDNNF} can be viewed as a decision diagram with additional structured decomposable $\land$-nodes. Conditioning is easy on branching programs: replace each decision node labeled by a conditioned variable with the child consistent with the conditioning term. The extra $\land$-nodes cause no problem: recurse through them, and structured decomposability is preserved because conditioning only removes variables from subcircuits. \citet{Wegener_2000} gives this argument for EXOR-OBDDs in Theorem 10.5.1.
\end{description}
\end{transformationclaim}

\begin{transformationclaim}
\claimlanguage{\langfam{dec-SDNNF}{T}}
\operation{\CD}
\status{\poly}
\begin{description}
For each fixed vtree $T$, a \langfam{dec-SDNNF}{T} can be viewed as a decision diagram with additional decomposable $\land$-nodes structured by $T$. Conditioning is easy on branching programs: replace each decision node labeled by a conditioned variable with the child consistent with the conditioning term. The extra $\land$-nodes cause no problem: recurse through them, and the fixed vtree is preserved because conditioning only removes variables from subcircuits. \citet{Wegener_2000} gives this argument for EXOR-OBDDs in Theorem 10.5.1.
\end{description}
\end{transformationclaim}

\begin{transformationclaim}
\claimlanguage{\langfam{cSDD}{T}}
\operation{\NOTC}
\status{\poly}
\begin{description}
\citet{VanDenBroeck_2015} proved that reduced \langfam{SDD}{T}{s} support negation in polynomial time, by passing negation through the subs.
\end{description}
\end{transformationclaim}

\begin{transformationclaim}
\claimlanguage{\langfam{cSDD}{T}}
\operation{\ANDBC}
\status{\nopolyunknownquasi}
\begin{description}
\citet{VanDenBroeck_2015} proved that combining two \langfam{cSDD}{T} formulas can force exponential \langfam{cSDD}{T} size.
\end{description}
\end{transformationclaim}

\begin{transformationclaim}
\claimlanguage{\langfam{cSDD}{T}}
\operation{\CD}
\status{\nopolyunknownquasi}
\begin{description}
\citet{VanDenBroeck_2015} proved that conditioning a \langfam{cSDD}{T} formula on a literal can force exponential \langfam{cSDD}{T} size.
\end{description}
\end{transformationclaim}

\begin{transformationclaim}
\claimlanguage{\langfam{cSDD}{T}}
\operation{\SFO}
\status{\nopolyunknownquasi}
\begin{description}
\citet{VanDenBroeck_2015} proved that forgetting one variable from a \langfam{cSDD}{T} formula can force exponential \langfam{cSDD}{T} size.
\end{description}
\end{transformationclaim}

\begin{transformationclaim}
\claimlanguage{\langfam{nOBDD}{<}}
\operation{\ORC}
\status{\poly}
\begin{description}
Disjunction of \langfam{nOBDD}{<}{s} is obtained by adding a top-level nondeterministic $\lor$-node whose children are the input formulas. Since all inputs respect the same fixed order $<$, the result is still a \langfam{nOBDD}{<} formula. \citep{TCZ_Forthcoming}
\end{description}
\end{transformationclaim}

\begin{transformationclaim}
\claimlanguage{\langfam{nOBDD}{<}}
\operation{\ANDBC}
\status{\poly}
\begin{description}
The usual product algorithm for conjoining \langfam{OBDD}{<}{s} respecting the same variable order works for \langfam{nOBDD}{<}{s}. Run the two diagrams in parallel according to the fixed order $<$, and accept exactly when both components accept. The resulting diagram still respects $<$. \citet{Wegener_2000} gives the analogous construction for EXOR-OBDDs in Theorem 10.5.1.
\end{description}
\end{transformationclaim}

\begin{transformationclaim}
\claimlanguage{\langfam{nOBDD}{<}}
\operation{\CD}
\status{\poly}
\begin{description}
Conditioning is easy on branching programs: replace each decision node labeled by a conditioned variable with the child consistent with the conditioning term. The fixed order $<$ is preserved. \citet{Wegener_2000} gives this argument for EXOR-OBDDs in Theorem 10.5.1.
\end{description}
\end{transformationclaim}

\begin{transformationclaim}
\claimlanguage{\langfam{nOBDD}{<}}
\operation{\FO}
\status{\poly}
\begin{description}
To forget a variable, replace each decision node labeled by that variable with a nondeterministic $\lor$-node over its two children. Repeating this for all forgotten variables preserves the fixed order $<$ and computes existential quantification. \citep{TCZ_Forthcoming}
\end{description}
\end{transformationclaim}

\begin{transformationclaim}
\claimlanguage{\langfam{uOBDD}{<}}
\operation{\CD}
\status{\poly}
\begin{description}
Conditioning is easy on branching programs: replace each decision node labeled by a conditioned variable with the child consistent with the conditioning term. The fixed order $<$ is preserved, and the operation cannot create two accepting paths for a remaining assignment. \citet{Wegener_2000} gives this argument for EXOR-OBDDs in Theorem 10.5.1.
\end{description}
\end{transformationclaim}

\begin{transformationclaim}
\claimlanguage{\langfam{uOBDD}{<}}
\operation{\ANDBC}
\status{\poly}
\begin{description}
The usual product algorithm for conjoining \langfam{OBDD}{<}{s} respecting the same variable order works for \langfam{uOBDD}{<}{s}. Run the two diagrams in parallel according to the fixed order $<$, and accept exactly when both components accept. The result is still unambiguous: for any satisfying assignment, each input has at most one accepting path, so the product has at most one accepting path. \citet{Wegener_2000} gives the analogous construction for EXOR-OBDDs in Theorem 10.5.1.
\end{description}
\end{transformationclaim}
\bibliographystyle{plainnat}
\bibliography{refs}
\end{document}
